\chapter{岁计与人口}

简十四实为全书账目考订之枢纽，其首行仅十六字：「计口三十一万二千，实居十一万九千。」两项数字之间相差十九万三千，其中前十六万二千人系在建造期内消耗损减：依据第 117 全宗第 2187 号卷宗及俄方先期复原之工程时序，该城池原系按三十一万二千人之额定容量设计动工，然迨至竣工之日，该族群总口数已跌落至十五万。是故，此城自正式启用之初便注定处于半数空置之境地；其后两万年间由十五万递减至十二万，再减至撤离前夕之十一万九千，不过是在既定空置格局之上继续向下收缩而已。刻写者将「计口」与「实居」并排刊刻于同一枚残片之上，其意图本身即足以证明，当时的书吏深知二者落差悬殊，且早已将此种永不满员之状态视为日常行政的前提。

\section{谁吃饭}

在推进具体口粮核算之前，首须解决一个先决前提：卷二全卷十片之中多达六片涉及粮食配给，载明「粟支九年」「月给粟，人日二升」「定额：人日三合」以及「赏以粟，罚以役」诸条，整套账目显系预设城内存在一批必须仰赖谷物维生之人口；然而，依照第一章所考订之月之民生理特征，月之民本身并不进食。这一判断虽未见残简以直白语言道破，却可自文献分布中严密反推：卷二专记口粮，卷一简四专记壁中气路与水路；假使居住环内之城主同样依赖谷物果腹，则卷二十片乃至全书三十七片之中必当存有配给居者正餐之规程，然而遍检全部残简，竟无一字涉及月之民自身之饮食消费。唯一表面相涉之简十五「月给粟，人日二升」，若联系前后片文义读之，其配给对象实指外部劳役群体而非环内城主。

若月之民确不食粟，卷二六片所反复核算之庞大粮秣究竟供予何人，逻辑上无非三种可能：其一，供予承担捣药苦役之月兔群体，此时简十五之主语即指「捣者」及杂役；其二，供予居于垣外小舍之三百二十名外工及往来地面之人员；其三，供予城中少数尚未完成体质改造之过渡人口。上述三种可能彼此并不排斥，笔者倾向于前两种解释并行，亦即卷二之粟与口粮系用以支应劳役与外工队伍之成本，而非城主群体之日常消费。此项重读将彻底改变卷二数字之历史分量：倘若粮秣仅发给做工一方，则简十二「粟支九年」所指者乃是维持一支庞大劳动队伍运转九年之周期，而非供养整座居住城池九年之储备，同一条记载的工程含义亦随之截然有别。

\section{三册分掌}

简十一明确交代了城内人口之统计与分管体制：「匠籍别为一册。工籍别为一册。居者不入册，以口计。册各有吏，不相知。」依简十所列总数「岁计口，凡三十二万一千」计之，全城名义人口遂被切割为三个互不统属之板块：入匠籍者六万，入工籍者十一万，另余无籍而仅按人头统计之居者十五万一千。前两项共计十七万人各立专册，第三项占比最大之居者群体反而全无册籍可稽。此项编制格局直接决定了本书所能复原的账目边界：凡有详确分项记录可循之数字，悉数出自那十七万有籍人员之名册，至于十五万一千居者，除了总数之外几无细账可查。尤为关键者在于「册各有吏，不相知」六字：三类人员既由互不相通之书吏分掌，城中便不存在任何一个能够总览十七万有籍者全貌的单一行政节点，至于不入册之居者更被彻底排除在名籍追踪体系之外。

\section{总差额}

将简十四额定人口三十一万二千与实居十一万九千相减，全期总差额为十九万三千人。在此十九万三千人之中，简十九所载兔乱事变死者仅占二万二千，余下之十七万一千人皆为建城及漫长运转期间逐步流失之数。在一个丧失生育能力之封闭群体内，人口走势只能呈现单向递减，每死亡一人即永久少一人，而偶发性新入者远不足以抵消损耗；简十三特设「新入者，先居外舍三月」之隔离程序，反证外来补充实属极少数特例，全书亦从未记载任何大批量入城之例。由此，全城之人口总账实可清晰划分为三个阶段：

\begin{center}
\small
\begin{tabular}{@{}llr@{}}
\toprule
段 & 内容 & 人数\\
\midrule
第一段 & 建城至完工，自然减员 & 三十一万二千降至十五万\\
第二段 & 六批迁移 & 十五万降至十二万\\
第三段 & 最后数年的零头与尾数 & 十二万降至十一万九千\\
\bottomrule
\end{tabular}
\end{center}

简十九之二万二千仅能解释第二段以后之局部骤减，而决不能将第一段以降两万年间之运转视作平缓无波。简十二与简十七所保留的两次局部危机实证，恰好揭示了城内资源吃紧时之真实运行逻辑。简十二末段记某年存粟仅支三月，当局遂依「减自上始」之序令居者日减为一升而工匠如故；按简十五「人日二升」计，一升正为半额，以十五万一千居者计之，此举一年所节存之粮秣约当总仓容之十四分之一。简十七末段则记十四年十一名捣者连日未足定额，各罚役三日而不改籍，直至其中一人再犯方予移籍。两案并观，足以窥见该行政体制之冷酷秩序：粮食短缺时优先削减最无议价权之居者定额，劳动欠额时则严格依循有限惩罚梯次，直至逾越三次上限方才动摇名籍，一切皆依既定规程运转，绝无情绪化之法外开恩。

\section{八千之数}

在全部数字核校中，第五章末已先提及之「八千」差额，乃是笔者在本卷内唯一无法妥善归位之数据。按第 117 全宗第 2187 号卷宗复原，设施完工时人口为十五万；其后分六批共迁离十二万，两万年间实减三万人；然简十九所录全部事变死者仅为二万二千，三万减二万二千，尚余八千人毫无下落。

关于此八千人之去向，逻辑上存在三种可能解释。其一，简十九所记之外尚存在未被入账之损耗：在两万年之时间跨度内，死亡缘由断不可能仅限于「兔乱」一宗，该事变之所以单独入简，或许仅因集中暴发之死伤需要行政处置与核销，而零散发生、无须专项善后之死亡则根本不予立档。其二，简十四「十一万九千」或系取整之数：残片该处恰有一格漫漶蚀去，缩微胶片本作「十一万九千六百」；倘若原数果为十一万九千六百，则差额将缩减至四千左右。\footnote{按简十四上下两处皆作带「百」字之三位数表述，且同片「计口」一栏尚存半笔刻痕；然原件此处已蚀，胶片本所作「六百」二字无从实勘，故上编校记未予径补，仅存此备考。}其三，十五万这一基数本身存在统计误差：该数系莫斯科档案馆第 117 全宗第 2187 号卷宗据另一批工程档案反推之复原值，而三十七片残简本身并无此数，倘若俄方所据原始工程统计略有出入，则此处差额便无从在本书内部校平。\footnote{第 117 全宗第 2187 号卷宗仅录完工时十五万之总数，未附分项明细账目，本书三十七片残简亦不载此数。}三种解释之中，笔者倾向第一种，然明言并无直接文献证据；之所以仍将此疑案详列于考释正文而不移入上编存疑，盖因涉及此数之原始文字俱清楚无讹，真正存在缺环者乃是史料背后的推论链条。

\section{旁证}

尽管无法终局定案，简十五与简十二合观所得之粮秣推算，仍可为第一种解释提供一项有力的间接旁证。简十五定额为「人日二升」，以此估算，十一万九千人日需谷物二千三百八十石，年需八十六万八千石；参以简十二「粟支九年」之数，九年总仓容当在七百八十一万石之巨。倘若城中粮仓确具如此庞大之储备规模，则在撤离前相当长之时期内，全城绝不至于面临饥荒危机；真正迫近临界极限而不可逆转者乃是空气储备。简二十七「急则气先绝，水次之，粟又次之」一语，已将三种维生物资之耗竭先后刻画得十分明晰，撤离行动必然触发于气路将绝而存粟尚丰之际，此亦与简二十四撤离装载时「粟尽、水尽、气尽」全数携走之记录全然相容，所携走者正系尚未耗尽之库底余粮。据此推之，那八千名未见于死伤账目之人群，绝少可能死于饥馑，而更可能属于「无须动用专项处置程序之常规死亡」，从而在账簿视野中自然隐没，此点正与前述第一种解释彼此呼应。

\section{章末自白}

在结束全卷核心数字之考辨之际，有必要就本书的数理校勘体例作一交代。凡三十七片内部可以交叉核验之数字，笔者一律逐笔核算；其确实无法弥合之处，皆如实悬置于正文及注释之中，绝不曲为弥缝。读者若对某一项折算比例或差额归因另有见解，完全可依据附录一所列全部原始编号与尺寸数据重行演算，原件既在，此项核算向任何人敞开。这一声明绝非学术谦辞，盖仅凭区区三十七片劫余残简所拼合而成之一本残账，在客观上本就绝无可能首尾凑平。

\chapter{撤离的经过}

简二十九与简三十集中交代了六批迁离之具体人数与计划属性，简二十至简二十八则逐项记录了撤离最终一年内所执行的一连串封闭与清退程序。本章将此两组材料并案排比，试图复原公元前二百二十一年前后那场撤离行动之真实面貌；就档案所呈示之程序纪律而言，那一年所发生者乃是一次严谨而冷酷的资产清点与设施退役，实非仓皇溃逃。

\section{六批}

\begin{table}[htbp]
\centering
\small
\caption{六批迁移的人数与年代}
\label{tab:six}
\begin{tabular}{@{}llr@{}}
\toprule
批次 & 时间（据第 2187 号卷宗） & 人数\\
\midrule
第一批 & 约前 421 年 & 6\,000\\
第二批 & 约前 380 年 & 8\,000\\
第三批 & 约前 340 年 & 10\,000\\
第四批 & 约前 300 年 & 13\,000\\
第五批 & 约前 260 年 & 18\,000\\
第六批 & 前 221 年 & 65\,000\\
\midrule
合计 & & 120\,000\\
\bottomrule
\end{tabular}
\tabsource{人数据本书简二十九；年份据第 117 全宗第 2187 号卷宗与《Лунная прародина》所记间隔推算，非三十七片所具列。}
\end{table}

需先说明者，本书所辑三十七片仅给出六批人数，并未附具具体公元年份；表\ref{tab:six}所列年代，系据第 117 全宗第 2187 号卷宗及《Лунная прародина》所录「每隔四十年前后」之周期记录，以第六批落地年份公元前二百二十一年为基准向前回溯推算而得，档案原文并无六栏年份明文，唯有各批人数与本书简二十九分毫不差。观察六批之时间跨度，每批之间相隔约四十年，自第一批至第六批首尾恰好历时二百年，这一时间长度绝非偶然巧合：它正与简二十七所记地热衰竭周期「二百年间，三十一分」完全重合。由此观之，所谓六批撤离，实为一段长达二百年的热源衰退曲线上的六个预设刻度，其间并不存在六次临时起意的独立决策，前五批始终以约四十年为固定步长逐批定额减载，呈高度规律化之有序撤出态势。

\section{最后一年}

简三十以十二字划分前后两段撤离之本质差异：「前五批，岁计之内。第六批，岁计之外。岁计之内者，有次第；岁计之外者，无次第。」前五批既在常规年度预算与计划之内按部就班推进，所遗留者便仅余冰冷之汇总数字；反观第六批因跌出既定计划之外而丧失常规次第，反倒迫使书吏在最后一年里仓促刻下大量过程性执行记录。这构成了本卷史料学上一桩颇具反讽意味的事实：恰因那一年原有秩序陡然溃坏，后世学者方得以窥见其溃坏之实际序列。若将简二十至简三十一依工程与行政逻辑重新编排，最后一年之运作经过大致可还原如下。

首道程序为「会计」清点。简二十载「元年，会计。在籍十一万九千。可去者十一万九千。留者无」，此「会计」一词原非凭空而来，其触发扳机正刻于简二十七末段：末年之际，壁中气路三度断裂而又三度勉强接复，司事之吏逐级上禀，主者遂定下「再断，则会计」之终局指令；待第四次气路断裂发生，全城人口清点随即全面启动。「三断三复」四字尤足深味：它表明该设施此前已历经多轮局部故障，屡次损毁而又屡次抢修维持，及至最后一年主者终于下定不再修复之决心。从工程后果观之，公元前二百二十一年之最终撤离实属环境系统彻底失效后之被迫止损，当年生态与维生网络全面崩溃，抽取污秽之终端段停止运转，环内污秽失控，事态已然超越修葺与否之工程讨论范畴。简二十将「在籍」「可去者」两数并列等同，继以「留者无」作结，复以「凡不去者，其名不入此册」作补足，从而在簿籍定义上径直将一切留守人员排除于统计范围之外。

继「会计」之后即为核算装载。简二十四载「装载：人十一万九千。粟，尽。水，尽。气，尽。器，择要者三千件。册，尽」，五类物资全数清库带走而不留余物。其中「册，尽」一项与后文退役清单第十七项「名，削」显相抵牾，构成卷三两处难以弥合的硬性矛盾之一。装载既定，随即核发登船凭证。简二十五载「凡去者，各给符一枚。符不还。符上刻其名。无名者不得登。符失者，不得补」，以具名之符作为登船唯一凭据，此复与第十七项「名，削」形成第二处硬性矛盾。笔者在上编校记中曾提出一种折衷推想，即先给符登船、后在底册削名，然三十七片中绝无任何一片可证实此种前后时序，故该解释终究仅止于无据之揣测。最后阶段，方才进入十九项退役清单之逐条处置。

\section{十九项的顺序}

简二十一至简二十二所列退役清单的条目序号，在纸面上未必等同于施工队列的实际作业次序，然而从工程依赖关系审视，该设施在技术上几乎只能依循此一脉络逐项关闭。第一项封居住环、第二项断气路、第三项断水路必系同步推进，盖气路一经切断，居住环即刻丧失维持能力，水源亦无从独存。其后第四项封库、第五项焚籍、第六项毁器，构成清场销毁之核心步骤。第七项拆除垂直运输电梯，将支柱于四十米高度截断、填土夯实、覆以三层石板并植草，且在截断处刻记「勿挖」二字；此条系全清单篇幅最长之一项，亦为唯一明确留有地面施工标记之条款。继之以第八项断界、第九项平外舍、第十项塞垣门，至此，作为居住与交通枢纽之城池本体已彻底丧失运转机能。

清单后半段九项（第十一至第十九项）所处置者，则皆为无法直接炸毁或另有深意之体制要素：封山（十一）、熄炉（十二）、弃药（十三）、收册（十四）、移界石（十五）、除标（十六）、削名（十七）、减口（十八），直至最后一道工序「门，闭」（十九）。第十九项「门闭」构成了全书三十七片最后一个具体动作，简三十一以罕见的白描笔法记下当日景象：「门闭之日，无送者，无哭声。闭者七人，回首一次。道上无迹，有迹者风。」七人闭门，既无送别之人，亦无哭泣之声，关门者仅仅回首一次；句中不具主语，亦未交代此七人此后去向。参以第九项已将外舍夷平、第十项已将垣门堵塞，倘若此七人系自垣内撤出，则退路早已断绝；倘若系于垣外最终闭门，则此七人从此不复入内，全部官方记录至此戛然而止。

\section{外工三百二十人}

简二十八之存留，直接撕开了最后一年清点账目中的一道裂口：「外工三百二十人，不入装载之数。留。留者，给三月之粮。」此三百二十名外工既被排除在十一万九千装载名额之外而遭留置原地，又获发三个月口粮，足证其性质属于短期留守而非即刻弃杀；然而三月粮尽之后何去何从，连同此项留守究竟发生在第九项「平外舍」之前抑或之后，残简均无只字交代。倘若先平外舍而后人员退走，此三百二十人势必无处容身乃至掩埋于废墟之中；倘若人员先散而后平舍，其去向亦复杳然无踪。更致命之矛盾在于，简二十明载「留者无」，简二十八却明载「留」，两片在编目号上仅隔八号，出自同一批移交材料且同属一九六三年编目序列。笔者不打算强行调和此两种相互排斥之记载：它们构成了本卷最为赤裸的一处文本冲突，而冲突本身即是最诚实的历史证词，足证在公元前二百二十一年那个秩序溃散的撤离年份里，负责造册的书吏之间同样未能将账目对齐。

\section{去者}

第六批六万五千人落地之后散落于豫西、关中一带之具体分布，属于地面着陆档案与第 117 全宗第 2187 号卷宗之复原成果，本书作为辑佚专著不予详述，仅能自残简内部补出两项值得审思的微观细节。其一即简三十一「无送者」三字：六万五千之庞大迁徙队伍启程离城，按常理绝无可能全城寂然无送，除非在当时之政治或体制语境下，这是一场没有任何留守机构或主体愿意公开承认的离去。其二见于卷四简三十五「□有十二，不可尽。尽者三，余九亡」：此「十二」之数恰好与相关工程档案中运载器乘员舱总数相合，然通篇残简并无任何一句能够坐实二者之直接对应关系，故本书严格恪守孤证不立之原则，不据此衍生立论。

\section{撤离之后}

自公元前二百二十一年末门闭合之后，该城之连续行政记录即告彻底中断，直至三百七十二年之后（约当公元一百五十一年），方始再见外来者之题记痕迹：简三十二载「后三百七十二年，有客至。门不可启，壁上题名而去」。然而另一方面，简三十三却复自称「此舟也。舟已不行。吾辈守之三世矣。守者无禄，亦无代」。倘若简三十三所言非伪，则意味着在那三百七十二年之静默期内实有人留守而无正式文书留存，且简中「吾辈」不言来历，三世跨越数百年亦与常人寿限悬殊。有守者而无公文记录，此一反常现象本身即具有制度史意义：它再次印证了该城之书写体系绝非自动运行，而是高度依附于行政建制；一旦原有行政建制解体，即便尚有残余守者存活，成系统的官方记录亦随之永久停摆。

\section{章末自白}

在本章收束之际，必须就上述最后一年撤离过程之写法作一诚实剖白。以上所排定的封闭、给符与退役时序，全系笔者个人依据动作逻辑与工程常识，将散乱无序之十二片残简强行串联而成；原件之上既无编年，亦无月份标示，所有「先」「后」次序纯粹出于技术可能性之推演，换作另一位研究者来排比，完全可能得出另一套截然不同的时序。凡涉及先后步骤之判断，读者皆可据技术理路予以推翻。不过，全章之中尚有两项结论并不依附于任何主观排序，因而具有较为坚实的文献可靠性：其一为十九项退役举措依封、断、毁、弃、移五类性质所作的量化分类，其二为简二十「留者无」与简二十八「外工留」之间不可调和的制度冲突。前者足以确证此城系经由系统化处置而主动关闭，实非无序弃置；后者则足以证实执行处置的各部门之间，至终亦未曾把账目理清。

\chapter{名与籍}

第七章既已考辨城池撤离之终局，本编最后一章将视线收束回一个贯穿全部三十七片之核心制度枢纽，即「名」与「籍」之关系。全书所录户籍、器记、符牌与焚籍条款皆由此衍生，下文将依次比勘四处关键经文，以检验此套名籍体制如何支撑整座城池之日常运作。

\section{名在册上}

在展开逐片比勘之前，必须先行明确一项基本界定：在这座城池的行政语汇中，所谓「名」，本质上指代名册上的一条登录条目，而非现代哲学或文学意义上之人格称谓。简十九「逃亡者，籍除，名削」、简二十二第十七项「名，削」以及简二十五「无名者不得登」「符上刻其名」，此三条所言实为同一套管理机制：入环则给籍，逃亡则除籍削名，无名籍者不得登船，可见「名」之有无直接等价于名册上是否存有该条登记。若循此界定反观卷一最难解之简九，诸多晦涩之处便豁然而通：简九将「呼山名」与「越界」判处同罪，表面视之二事风马牛不相及，落入册籍体制之下却属同一性质之违制，盖越界乃是擅自篡改自身在册之坐标位置，而直呼山名则是擅自替册外之物强行立条入籍。在城中，立册、除籍、移籍皆系专职吏员之排他性权限，普通居者或役夫妄行越界与呼名，均构成对册籍专断权之僭越，故依法判为同罪；简九实乃一条严厉的户籍治安条款，绝非一条关于言语或存在的形而上学格言。三十七片之中「名」字频见于四个关键场域，即简十三「旧器有记，记其旧主」、简十九「逃亡者，籍除，名削」、简二十二第十七项「名，削」以及简二十五「无名者不得登」「符上刻其名」，本章以下即依次将此四处并置比勘，以检验其背后是否遵循同一套制度逻辑。

\section{记名}

首先考察简十三之「记名」条文：「新入者，先居外舍三月，乃入环。入环之日，给籍，给衣，给器。器不继，以旧器补。旧器有记，记其旧主。」新入环者所领受之籍、衣、器三物悉数由公府统一配发，若遇器物不继而改以旧器递补，旧器之上仍保留原持有者之标记，而不予刮除。此处必须审慎辨析「记」与「名」之用字差异，尤须留意全书三十七片中绝无任何一人留有私名这一铁的事实：假使旧器之上所刻果真为私人姓名，历经两万年流转，这些名字理应随实物残留至今，而考古编目所得竟无一字。因此，简十三所谓「记其旧主」，所刻极可能实非个人姓字，而系籍号、环号一类的行政编码；倘若此读法成立，则简九之禁呼私名与简十三之留存旧记便各不相悖，盖所禁者为「名」，所记者为「籍」。当然，倘若此读法不成立，简十三便将直接构成第三章所述净秽清除原则之反例；二者之间笔者倾向于前者，然同样声明并无直接实物佐证。退一步言，即便容许旧器留存追溯标记，在一个名册尚且岁一焚毁之体系内，器物记号供物资流转追溯之用，名册焚毁供体制定期清空之用，二者分管不同行政维度，原可并行而不相抵触。

\section{削名}

次须比勘两处形态相同而对象迥异之「名，削」：简十九载「逃亡者，籍除，名削」，简二十二第十七项则载「名，削」。二者同用「削」字，所针对之尺度却大相径庭：简十九所削者乃单个逃亡者之个体记录，属于违纪惩戒；简二十二所削者乃撤离前夕全城之总体名籍，属于彻底清场。在同一套官僚制度之中，「名」既充当追查核验之唯一凭据，同时又充当最终注销清算之靶向对象；表面观之似属自相矛盾，实则恰为一事之两面：正因名字预先登载于册籍之上，方具备可追查之行政效力；亦正因名字被固化为册上条目，在设施终结之际方能一并予以技术性削去。倘若名字本不在册，则平日之追查与末日之清算均无从谈起。

\section{无名}

全卷关于名籍最棘手之悖论，集中爆发于简二十五与简二十二第十七项之直接碰撞。简二十五规定「凡去者，各给符一枚。符不还。符上刻其名。无名者不得登。符失者，不得补」，明定登船者必先领符、符上必刻名、无名则不得登船；然而紧随其后的退役清单第十七项偏偏赫然列着「名，削」。究竟系先发符登船、后削底册之名，抑或先削城中名籍、再据别籍发符，全部三十七片内并无任何一片足以裁定这一执行先后；笔者在上编简二十五校记中虽曾给出一种时序调和之假说，然彼语纯属无据之推想，绝不能充作定论。单就现存文本所能证实之极限而言，若要令此二条同时成立而不致逻辑崩溃，唯一的出路在于承认「名」在撤离那一年内已分裂为两套独立载体：其一为刻于登船符上、随逃生者个人携行带走之「名」，其二为留存于城内底册、供留守清场人员最终抹除之「名」。唯有作此区分，方能解释为何「名，削」被排在退役清单接近末尾之第十七项，盖其本即城池销毁程序之倒数第三道收尾手续，其抹除对象乃留驻原地之行政底册，而非正待登船撤离之活人。

\section{一个空缺}

至此，全部考释不得不重新回到自第一章起便始终悬而未决的那项最刺眼的档案事实：三十七片残简之中，竟没有留下任何一个具体的个人名字。四片反复围绕「名」字展开规程与惩戒之残简，详尽讨论了记名、削名、移籍与无名不得登船，却偏偏没有任何一片写下过哪怕一个具体人名，甚至连那名因「欲见山」而敢于违禁越界的外工，亦只在档案中存留了三字理由，未尝留下一姓一名。倘若第一章末尾所谓该体系「对名的态度并不统一」可以成立，那么这种彻底的人名空缺便绝非材料散佚所致之偶然，而是制度运作自身制造出的必然结果：一方面在日常中强制将人绑定于籍册编码，另一方面在年终岁焚与最终撤离时复将名籍彻底抹除；在这两万年间「岁一易、岁一焚」的持续销毁机制过滤之下，侥幸残存至今者仅此未及付之一炬之数十片断简。本书作为辑佚考释所能提供之全部实证材料到此为止；至于后世读者必然会追问的「那么那些居于环内之人各自究竟叫作什么」，在这三十七片残简之中，找不到一个字的答案。

\chapter*{结\quad 语}
\phantomsection\addcontentsline{toc}{chapter}{结\quad 语}
\markboth{结语}{结语}

导论是从一九六三年莫斯科那间旧库房开始的，全书的收束也回到同一间库房。这三十七片残简在那里静置了六十年，其间经历过三次编目、一次馆藏转移，以及两次险些被当作无归属杂物清理出库的危机：第一次是在一九七〇年代馆舍扩建之时，一批「无归属」箱件被列入清退清单；第二次是在一九九一年机构改组之际，无人接手之冷门卷宗被划归暂存待议。两次清退最终皆未执行，原非因为库房里站着哪一位慧眼识珠的守护者挺身阻拦，纯粹是因为彼时经办之人事务繁忙、来不及逐箱清理。此种存留方式固然缺乏体面，却最接近档案史实的真相：这批铭刻着月面城市终局的薄片之所以能够侥幸传世，仅仅因为在每一次档案清理的序列里，它们恰好被堆放在最靠后的位置。

我在导论里已经写明，这一卷三十七片的编排次序出自我的主观交叉比对，换作另一位研究者必然会排定出另一副面貌；释文中约有一成缺字出于我的推补，读者亦完全有权持保留乃至推翻。行文至此，还须坦白补充一句：不只我的工作建立在推论之上，这整部书得以问世的物质前提，同样建立在那两次「尚未及清理」的官僚拖延之上；倘若当年清退人员的动作哪怕快上一步，世间便不会存在这一卷书。就内容本身而言，这三十七片中没有任何一位英雄人物的踪迹，所能见到的唯有一年一检的管壁、岁一焚毁的名册、日三合的劳役定额，以及最后一年里冷冰冰的十九项退役清单；全书唯一近似故事的片段，仅仅是一名不知姓氏的外工因想望一眼界外之山而越界受罚，留下「欲见山」三个字。如今我已将这本书写完；至于它日后能否在学术书架上存留下来，其缘由大概也与那三十七片残简别无二致：未必因为终于有人认得了它的价值，而多半只因它在清退名单上恰好排在最后。

\begin{flushright}
{\small 巴兰尼科夫}
\end{flushright}
